% 第4章 4.2.2：(a) Cantor 函数 c（8 层端点折线，端点值精确），在挖去的区间 (1/3,2/3)、(1/9,2/9)、(7/9,8/9) 上分别取 1/2、1/4、3/4；
% (b) 测试 phi 只在 (1/9,2/9) 与 (7/9,8/9) 内变化（阴影），-∫ c phi' = -(1/4)(1)+(3/4)(1)=1/2。数据由同目录 generate.py 计算。
\begin{tikzpicture}[font=\figurefont\footnotesize]
\pgfplotsset{every axis/.style={width=52mm,height=40mm,axis lines=left,
  x axis line style={draw=FigureStroke,line width=.6pt},y axis line style={draw=FigureStroke,line width=.6pt},
  tick style={draw=FigureStroke},tick label style={font=\figurefont\scriptsize,text=FigureInk},
  xlabel={$x$},xlabel style={at={(axis description cs:1,0)},anchor=west,font=\figurefont\scriptsize},
  title style={font=\figurefont\footnotesize},xmin=0,xmax=1.02,ymin=0,ymax=1.08,
  xtick={0,0.3333,0.6667,1},xticklabels={$0$,$\frac13$,$\frac23$,$1$},ytick={0,0.25,0.5,0.75,1},
  yticklabels={$0$,$\frac14$,$\frac12$,$\frac34$,$1$}}}
\def\dat{figures/02-abstract-mathematics/03-function-spaces/tikz/c4-cantor}
\begin{axis}[name=A,title={(a) Cantor函数$c$}]
\addplot[draw=FigureBlue,line width=1.1pt] table[x=x,y=y] {\dat/cantor.dat};
\end{axis}
\begin{axis}[at={(A.east)},anchor=west,xshift=12mm,title={(b) 只在两处平台上变化的测试},
  xtick={0,0.1111,0.2222,0.7778,0.8889,1},xticklabels={$0$,$\frac19$,$\frac29$,$\frac79$,$\frac89$,$1$},
  legend style={draw=none,fill=none,font=\figurefont\scriptsize,at={(0.5,-0.28)},anchor=north,legend columns=2}]
\fill[FigureYellowFill] (axis cs:0.1111,0) rectangle (axis cs:0.2222,1.08);
\fill[FigureYellowFill] (axis cs:0.7778,0) rectangle (axis cs:0.8889,1.08);
\addplot[draw=FigureBlue,line width=1.1pt] table[x=x,y=y] {\dat/cantor.dat};
\addplot[draw=FigureGreen,line width=1.1pt,dashed] table[x=x,y=y] {\dat/test.dat};
\legend{$c$,$\varphi$}
\end{axis}
\end{tikzpicture}
